% -----------------------------------------------------------------------------
% This file is automatically generated.
% Do not edit manually.
% Source: notebooks/support/hypothesis-testing/support.Rmd
% -----------------------------------------------------------------------------


\noindent If we draw a sample of size $n = 1,100$ from the \ttblue{accounts\\_receivable} population, the largest sampling units have an almost 100\% inclusion probability and with a selection method such as \ttblue{random} are likely to be hit more than once.

\noindent For example, the data file \ttblue{accounts\\_receivable}, that comes with the \ttblue{FSaudit} package, has the following six largest amounts:

\par\addvspace{\topsep}
\begingroup
\fvset{listparameters={%
  \setlength{\topsep}{0pt}%
  \setlength{\partopsep}{0pt}%
  \setlength{\parsep}{0pt}%
  \setlength{\itemsep}{0pt}%
}}
\begin{Verbatim}[breaklines=true,formatcom=\color{ada_blue}]
accounts_receivable[order(-accounts_receivable$amount), ][1:6, 1:3]
\end{Verbatim}
\begin{Verbatim}[breaklines=true,formatcom=\color{ada_light_blue}]
# A tibble: 6 × 3
  debtor   invoice amount
   <int>     <int>  <dbl>
1  47013 201719763 12050.
2  91415 201734940 11981.
3  33149 201706806 11586.
4  21704 201739548 10795.
5  39608 201726290 10520
6  65215 201723585 10377.
\end{Verbatim}
\endgroup
\par\addvspace{\topsep}

\noindent We set up a new \ttblue{mus\\_obj} and use a materiality of \ttblue{pm = 36730} to arrive at a sample size of 1100.
